% !TEX program = latexmk
\documentclass[aspectratio=169]{beamer}

% ---------- Theme ----------
\usetheme{Madrid}
\usecolortheme{default}

% ---------- Packages ----------
% LuaLaTeX handles Unicode natively — no inputenc/fontenc needed
\usepackage{fontspec}
\usepackage{kotex}          % Korean support (uses luatexko under the hood)
\usepackage{amsmath, amssymb}
\usepackage{graphicx}
\usepackage{booktabs}       % nicer tables
\usepackage{hyperref}

% ---------- Fonts ----------
% kotex picks a default Korean font automatically.
% Uncomment and change to a specific font if the default is unavailable:
% \setmainfont{Noto Serif}
% \setsansfont{Noto Sans}
% \setmainhangulfont{Noto Serif CJK KR}
% \setsanshangulfont{Noto Sans CJK KR}

% Korean italic workaround:
% CJK fonts have no italic variant, so \textit{} silently does nothing.
% \hangulit applies fontspec's FakeSlant (geometric slant) to simulate italics.
\newcommand{\hangulit}[1]{{\addhangulfontfeatures{FakeSlant=0.167}#1}}

% ---------- Metadata ----------
\title[Short Title]{Full Presentation Title}
\subtitle{Subtitle (optional)}
\author[Author]{Author Name}
\institute[Org]{Organization / Department}
\date{\today}

% ---------- Document ----------
\begin{document}

\begin{frame}
    \titlepage
\end{frame}

\begin{frame}{Outline}
    \tableofcontents
\end{frame}

% ==================================================
\section{Introduction}
% ==================================================

\begin{frame}{Motivation}
    \begin{itemize}
        \item First point
        \item Second point
        \item Third point
    \end{itemize}
\end{frame}

% ==================================================
\section{Background}
% ==================================================

\begin{frame}{Background}
    \begin{block}{Key Concept}
        Place a definition or important concept here~\cite{book:571187}.
    \end{block}

    \begin{alertblock}{Warning}
        Use alert blocks to draw attention.
    \end{alertblock}
\end{frame}

% ==================================================
\section{Results}
% ==================================================

\begin{frame}{Example Table}
    \begin{table}
        \centering
        \begin{tabular}{lcc}
            \toprule
            Method & Metric A & Metric B \\
            \midrule
            Baseline & 0.80 & 0.75 \\
            Proposed & \textbf{0.92} & \textbf{0.88} \\
            \bottomrule
        \end{tabular}
        \caption{Comparison of methods.}
    \end{table}
\end{frame}

% ==================================================
\section{한국어 테스트}
% ==================================================

\begin{frame}{한국어 (kotex) 테스트}
    \begin{block}{kotex 동작 확인}
        이 슬라이드는 \texttt{kotex} 패키지가 정상적으로 동작하는지 확인합니다.
    \end{block}

    \begin{itemize}
        \item 가나다라마바사 — 기본 한글 자모
        \item \textbf{굵은 글씨}와 \hangulit{기울임꼴} 지원
        \item 숫자 및 혼용: 2024년 1월 1일
        \item 영문 혼용: LuaLaTeX으로 컴파일된 한국어 슬라이드
    \end{itemize}

    \begin{alertblock}{주의사항}
        한글 폰트가 시스템에 설치되어 있어야 합니다.
        기본값은 \texttt{kotex}가 자동으로 선택합니다.
    \end{alertblock}
\end{frame}

% ==================================================
\section{Conclusion}
% ==================================================

\begin{frame}{Conclusion}
    \begin{enumerate}
        \item Summary point one.
        \item Summary point two.
        \item Future work direction.
    \end{enumerate}
\end{frame}

\begin{frame}[allowframebreaks]{References}
    \bibliographystyle{plain}
    \bibliography{refs}
\end{frame}

\end{document}
